\chapter{SYSTEM ARCHITECTURE AND IMPLEMENTATION}

\section{Compiler Pipeline}
The implemented system follows a four-stage pipeline. Source code is first
converted into tokens by the lexer. The parser then constructs an AST from the
token stream. After parsing, the type checker validates and infers types across
the AST. Finally, the interpreter executes the checked program and emits output
line by line.

The full processing pipeline is:

\begin{center}
\texttt{Source Code $\rightarrow$ Tokens $\rightarrow$ AST $\rightarrow$ Type Check $\rightarrow$ Execution Output}
\end{center}

This design keeps syntax handling independent from semantic checking and
execution, which simplifies testing and debugging.

\section{Project Structure}
The main source files and their responsibilities are summarised below.

\begin{table}[H]
\centering
\caption{Main Source Files}
\begin{tabular}{ll}
\toprule
File & Responsibility \\
\midrule
\texttt{main.py} & Command-line runner for sample code or a script file \\
\texttt{ui.py} & Desktop UI for editing and running a script \\
\texttt{components/tokens.py} & Token classes and token categories \\
\texttt{components/lexica.py} & Lexical analysis \\
\texttt{components/ast\_nodes.py} & AST node definitions \\
\texttt{components/parser.py} & Recursive-descent parser \\
\texttt{components/ast\_printer.py} & AST renderer for debugging/reporting \\
\texttt{components/type\_checker.py} & Static type checker and inference \\
\texttt{components/interpreter.py} & Tree-walking interpreter \\
\texttt{components/pipeline.py} & Shared pipeline for CLI, UI, and tests \\
\bottomrule
\end{tabular}
\end{table}

\section{Technology Stack}
The project uses Python and standard-library tools only.

\begin{table}[H]
\centering
\caption{Technology Stack}
\begin{tabular}{ll}
\toprule
Tool & Purpose \\
\midrule
Python 3 & Core implementation language \\
\texttt{dataclasses} & AST and runtime helper structures \\
\texttt{tkinter} & Desktop graphical interface \\
\texttt{unittest} & Automated testing \\
\bottomrule
\end{tabular}
\end{table}

\section{Lexical Analyser}
The lexer reads the source character by character and emits tokens for
identifiers, literals, keywords, operators, and delimiters. It preserves line
and column information in each token to support precise diagnostics. Invalid
characters and unterminated strings are reported immediately with source
positions.

Two-character operators (\texttt{==} and \texttt{!=}) are resolved with a
single lookahead step before single-character operators are matched. Identifiers
are scanned greedily and then checked against the \texttt{KEYWORDS} dictionary;
unrecognised names become \texttt{IDENTIFIER} tokens. Number scanning detects
the optional decimal point and lookahead digit that distinguish a
\texttt{FLOAT\_LITERAL} from an \texttt{INT\_LITERAL}.

\subsection{Implemented Lexer and Symbol Table Scope}
The lexer and symbol table work covers the full foundational scope assigned for
this component. Concretely, this stage produces the token stream consumed by the
parser and provides the symbol storage consulted by the type checker:

\begin{itemize}
    \item \texttt{TokenType} enum covering all four literal types, \texttt{IDENTIFIER}, six keywords (\texttt{if}, \texttt{else}, \texttt{while}, \texttt{def}, \texttt{return}, \texttt{print}), four arithmetic operators, two comparison operators, the assignment operator, six delimiters, and \texttt{EOF},
    \item frozen \texttt{Token} dataclass carrying \texttt{token\_type}, \texttt{lexeme}, \texttt{line}, and \texttt{column} for precise diagnostics,
    \item keyword mapping and two-character operator tokenisation in the \texttt{Lexer} class with source-positioned error messages,
    \item identifier scanning with keyword fallback; integer, float, Boolean (\texttt{true}/\texttt{false}), and string literal recognition,
    \item scoped \texttt{SymbolTable} with duplicate-symbol detection, parent-chain lookup, and child-scope creation,
    \item \texttt{LanguageType} enum (\texttt{Integer}, \texttt{Float}, \texttt{Boolean}, \texttt{String}, \texttt{Void}) used to label variables and function signatures,
    \item \texttt{VariableSymbol} and \texttt{FunctionSymbol} subclasses storing type, parameter list, and initialisation state,
    \item \texttt{format\_table()} method that renders the symbol table for the CLI and UI inferred-types display.
\end{itemize}

\section{Symbol Table}
The symbol table is implemented in \texttt{components/symbol\_table.py} and
serves as the shared store for both the type checker and the formatted output
displayed to the user.

\textbf{LanguageType enum.}\ The five language types are defined as an enum:
\texttt{Integer}, \texttt{Float}, \texttt{Boolean}, \texttt{String}, and
\texttt{Void}. The first four correspond to the language's primitive types.
\texttt{Void} is the inferred return type of functions that contain no
\texttt{return} statement, keeping the type system closed under function
definitions.

\textbf{Symbol class hierarchy.}\ The base \texttt{Symbol} class holds a
\texttt{name} and a \texttt{symbol\_type}. Two subclasses specialise it:
\texttt{VariableSymbol} adds an \texttt{initialized} flag that tracks whether
the variable has received a value, and \texttt{FunctionSymbol} adds a
\texttt{parameters} list of \texttt{(name, LanguageType)} pairs representing
the inferred function signature.

\textbf{Scoped lookup.}\ Each \texttt{SymbolTable} instance holds a
\texttt{\_symbols} dictionary for its own scope and an optional \texttt{parent}
reference. \texttt{lookup()} walks the parent chain so that inner scopes can
read outer variables. \texttt{exists\_in\_current\_scope()} checks only the
local dictionary, which \texttt{define\_variable()} and
\texttt{define\_function()} use to raise a \texttt{SymbolTableError} on
duplicate definitions. A new child scope is obtained by calling
\texttt{child\_scope()}, which is used by the type checker when entering blocks
and function bodies.

\textbf{Formatted output.}\ \texttt{format\_table()} renders the symbol table
as a fixed-width text table with columns for name, kind, type, initialisation
status, and parameters. This output appears as Stage~3 in the CLI pipeline and
in the inferred-types tab of the desktop UI.

\section{Parser}
The parser is implemented as a handwritten recursive-descent parser. It builds
AST nodes for assignments, print statements, returns, blocks, conditional
statements, while-loops, function definitions, and function calls. Operator
precedence is encoded structurally through separate parsing functions for
factors, terms, expressions, and boolean expressions.

Distinguishing an assignment statement from a bare expression statement is
handled with a one-token lookahead: when the current token is an
\texttt{IDENTIFIER} and the next token is \texttt{ASSIGN} (\texttt{=}), the
parser takes the assignment path; otherwise it falls through to the expression
path. This avoids backtracking and keeps the parser deterministic.

String literal values stored in \texttt{Literal} nodes have their surrounding
double-quote characters stripped during parsing, so \texttt{"hello"} produces a
\texttt{Literal} whose value is the plain Python string \texttt{hello}.

The \texttt{\_bool\_expr()} method handles the condition position in
\texttt{if} and \texttt{while} statements. It first checks for a bare
\texttt{BOOL\_LITERAL} and returns it immediately. Otherwise it parses an
expression and looks for a \texttt{==} or \texttt{!=} operator; if neither
is present the expression is returned as-is, and the type checker is responsible
for verifying that the result is \texttt{Boolean}.

\subsection{Implemented Parser and AST Scope}
The parser and AST work covers the full syntactic scope assigned for this
component. Concretely, this stage transforms the token stream into the tree
structure consumed by the type checker and interpreter:

\begin{itemize}
    \item recursive-descent \texttt{Parser} class producing a \texttt{Program} root node from a \texttt{list[Token]},
    \item \texttt{Assign} node: variable assignment with a target name and a value sub-tree,
    \item \texttt{BinaryOp} node: both arithmetic (\texttt{+}, \texttt{-}, \texttt{*}, \texttt{/}) and comparison (\texttt{==}, \texttt{!=}) operations encoded in the same node with an \texttt{op} string field,
    \item \texttt{Literal} node: integer, float, Boolean, and string values (string quotes stripped at parse time),
    \item \texttt{Identifier} node: variable and function name references,
    \item \texttt{If} node: condition, mandatory then-block, and an \texttt{Optional[Block]} else-block,
    \item \texttt{While} node: condition and body block,
    \item \texttt{FunctionDef} node: function name, list of parameter \emph{names} only (types are deferred to the type checker), and a body block,
    \item \texttt{FunctionCall} node: callee name and a list of argument sub-trees,
    \item \texttt{Print} node: single expression argument,
    \item \texttt{Return} node: single expression to return,
    \item \texttt{Block} node: brace-delimited sequence of statements used as function bodies, if-branches, and loop bodies,
    \item \texttt{Program} node: top-level list of statements,
    \item all nodes carry \texttt{line} and \texttt{column} metadata for source-positioned diagnostics.
\end{itemize}

\section{Abstract Syntax Tree}
The AST is the main intermediate representation shared by the checker and the
interpreter. It contains nodes for literals, identifiers, binary operations,
assignments, block statements, \texttt{if}, \texttt{while}, function
definitions, function calls, \texttt{return}, and the built-in
\texttt{print()} function. Each node stores line and column metadata.

\textbf{Node design decisions.}\ \texttt{FunctionDef} stores only the
parameter names, not their types; type information is not available at parse
time and is inferred by the type checker from the first call site.
\texttt{If.else\_block} is \texttt{Optional[Block]}: the \texttt{else} clause
is genuinely absent for single-branch conditionals, and the interpreter checks
for \texttt{None} rather than inserting an empty block. \texttt{BinaryOp}
serves dual duty — it represents both arithmetic and comparison operations,
distinguished at semantic time by the \texttt{op} string
(\texttt{+}/\texttt{-}/\texttt{*}/\texttt{/} versus \texttt{==}/\texttt{!=}).

\textbf{AST printer.}\ \texttt{components/ast\_printer.py} implements an
\texttt{ASTPrinter} class that walks the AST using a visitor dispatch: for each
node it looks up a method named \texttt{\_visit\_<NodeClassName>} and calls it.
The output is an indented text tree where each line names the node type and
its key fields (name, operator, value, argument count, etc.), indented four
spaces per level. This output appears as Stage~2 in the CLI pipeline and in
the AST tab of the desktop UI.

\section{Type Checker}
The type checker traverses the AST before execution. It uses the symbol table
to record variable and function information, infer types on first assignment,
validate arithmetic and comparison expressions, enforce Boolean conditions for
control flow, infer function parameter types from the first call, and infer
function return types from \texttt{return} statements. Programs that fail these
rules are rejected before interpretation.

Types are enforced entirely during this pass. The interpreter does not consult
the symbol table to validate each operation; it assumes the checked program and
uses native Python values.

\subsection{Implemented Type System Scope}
The implemented type system work covers the full semantic scope assigned for the
type system component. In presentation terms, this stage is responsible for
turning a syntactically valid AST into a semantically valid program that can be
executed with confidence:

\begin{itemize}
	\item static typing rules for arithmetic, comparison, assignment, control flow, and functions,
	\item type checking for expressions and statements,
	\item variable type inference from first assignment,
	\item function parameter inference from the first call,
	\item function return type inference from \texttt{return} statements,
	\item source-positioned type errors when a program violates a typing rule.
\end{itemize}

\section{Interpreter}
The interpreter is a tree-walking evaluator that executes the checked AST. It
uses nested runtime environments for variables and function calls. Assignments
store values in the current or nearest enclosing scope, \texttt{if} executes
only the selected branch, and \texttt{while} re-evaluates its condition until it
becomes false. Function calls create child environments, bind arguments by
value, and return through an internal return signal mechanism.

Bindings hold plain Python objects, not pairs of \texttt{(value, LanguageType)}.
The displayed language type on \texttt{print()} is derived from the object's
Python class (\texttt{int}, \texttt{float}, \texttt{bool}, \texttt{str}),
which matches the static rules because only type-checked programs reach this stage.

The built-in \texttt{print()} function outputs both value and runtime type, such
as \texttt{3 : Integer} or \texttt{"plc" : String}.

\subsection{Implemented Interpreter Scope}
The interpreter work covers all required runtime behaviours. In practical terms,
this stage is what gives the language observable behaviour from the user's point
of view:

\begin{itemize}
	\item assignment execution,
	\item \texttt{if} execution,
	\item \texttt{while} execution,
	\item function execution with value parameter passing,
	\item \texttt{print()} execution with immediate output,
	\item runtime integration through the shared pipeline used by the CLI, UI, and tests.
\end{itemize}

\section{Graphical User Interface}
The project includes a desktop UI built with \texttt{tkinter}. The interface
matches the requested layout: script input appears on the left, while output is
shown on the right. The right-hand side contains tabs for execution output,
tokens, AST, and inferred types. This allows the user to inspect each stage of
the language pipeline from one screen.
